\documentclass[10pt,a4paper]{amsart}
\usepackage[margin=19mm]{geometry}
\usepackage{amsmath,amssymb,mathtools}
\usepackage[scr=rsfs,cal=euler]{mathalfa}
\usepackage{xfrac}
\usepackage[unicode,hidelinks,bookmarks=false]{hyperref}
\newcommand{\ie}{i.e.\ }
\newcommand{\cf}{cf.\ }
\newcommand{\resp}{respectively\ }
\newcommand{\Subsection}{\subsection}
\providecommand{\nobreakdash}{\nobreak\hbox{-}}
\setlength{\emergencystretch}{2em}
\multlinegap0pt
\usepackage{bm}
\usepackage{bbm}
\usepackage{dsfont}
\usepackage{xspace}
\usepackage{cancel}
\usepackage{mathtools}
\usepackage{amsthm}
\makeatletter
\@ifundefined{theorem}{\newtheorem{theorem}{Theorem}}{}
\@ifundefined{claim}{\newtheorem{claim}[theorem]{Claim}}{}
\@ifundefined{lemma}{\newtheorem{lemma}[theorem]{Lemma}}{}
\makeatother
\title[Coloring Grids Avoiding Bicolored Paths]{Coloring Grids Avoiding Bicolored Paths}
\author{Derman Keskinkilic, Lale Ozkahya}
\date{Source: arXiv:2312.12919v7 (CC BY 4.0)}
\begin{document}
\maketitle

\noindent\emph{Source and licence.} Coloring Grids Avoiding Bicolored Paths. Version arXiv:2312.12919v7 (CC BY 4.0), distributed under the Creative Commons Attribution 4.0 International licence, as recorded in the arXiv licence metadata for this exact version and shown on the abstract page https://arxiv.org/abs/2312.12919v7. Full rights notice: licenses/arxiv-2312.12919-CC-BY-4.0.txt.\par\smallskip

\noindent\textit{Editorial scope.} Counted unit: the Lemma whose source label is \texttt{lem-case1} in the article (source0.tex lines 281--283, source label \texttt{lem-case1}), delivered together with the article's own complete original proof of it (source0.tex lines 284--321, one \texttt{proof} environment of 38 lines). No mathematics is written, repaired or re-derived: statement and proof are the article's own text line for line. Required context retained uncounted: lem-angle (lines 193--205). Every definition, display and label of those lines is retained unchanged. Omitted from this package and not redistributed: 1--192, 206--280, 322--386. Nothing inside a retained excerpt is omitted or altered: every retained body is delivered exactly as the article's lines are written, so no omission receipt of any kind accompanies this package. Citations: the retained excerpts carry no citation command at all, so no bibliography entry of the article is transcribed and no reference is redistributed. Reference adaptation: none was needed and none was made; every reference inside the delivered unit resolves inside it, so no unresolved reference or citation is produced. Printed slips kept literal: no printed word, symbol, hypothesis or number of the counted statement or of its proof was altered. Layout and class: the article's own class is \texttt{dmtcs-episciences} with packages including inputenc, subfigure, natbib, amsmath, amsfonts, amssymb, amsopn, amsthm; the wrapper is the lane's pinned \texttt{amsart} editorial wrapper at 10pt on A4, which restates the theorem-like environments the retained text needs and adds the article's own macro definitions that the retained text needs (none), with extra wrapper packages (bm, bbm, dsfont, xspace, cancel, mathtools). The delivered numbering is the unit's own. Assets: no retained span contains a figure environment, a graphics-inclusion command, a PDF-page inclusion, a table or a TikZ picture; no image file is copied into this package, the package declares no asset dependency and no image file is redistributed with it. Licence: version arXiv:2312.12919v7 of this article is distributed under the Creative Commons Attribution 4.0 International licence, as recorded in the arXiv licence metadata for this exact version and shown on the abstract page https://arxiv.org/abs/2312.12919v7. A full rights notice is delivered at licenses/arxiv-2312.12919-CC-BY-4.0.txt. Characters: every retained excerpt is pure ASCII, so the delivered engine, XeLaTeX with its default Latin Modern text font, reports no missing character.
\begin{lemma}\label{lem-angle}
For any $k\geq 5$, let $C$ be a (partial or peripheral) bicolored component in a 3-coloring $c: V(P_{k-2}\square P_{k-2})\to \{c_1,c_2,c_3\},$ and label the vertices along a partial walk $B_C$ as $v_1, v_2, \dots, v_r$.  
Then, the following hold:
\begin{enumerate}
    \item $r\geq 3$. 
    \item The angles between the edges $v_1v_2$ and $v_2v_3$ and between $v_{r-2}v_{r-1}$ and $v_{r-1}v_r$ are $90^{\circ}$.  
    \item For $i\leq r-2$, the angle between the edges $v_iv_{i+1}$ and $v_{i+1}v_{i+2}$ along $B_C$ is $90^{\circ}$ if and only if $i$ is odd. 
    \item $r$ is an odd integer.
    \item 
    Let $c(v_1)=c_1$ and $C$ have colors $c_1$ and $c_2.$ Then, 
    there is a $c_1$-$c_3$ colored connected subgraph containing the vertices $v_i$ and $u_{(i+1)/2}$, for odd $i$, $1\leq i\leq r-2,$ where $u_{(i+1)/2}$ is the $c_3$-colored vertex on the $C_4$ containing $\{v_i, v_{i+1},v_{i+2}\}.$
\end{enumerate}
\end{lemma}

\begin{lemma}\label{lem-case1}
    For any $k\geq 5$, if a 3-coloring of $P_{k-2}\square P_{k-2}$ has a peripheral bicolored component, then there is a bicolored $P_k$. 
\end{lemma}

\begin{proof}
If the 3-coloring of $P_{k-2}\square P_{k-2}$ has a peripheral bicolored component, then let $P$ be a maximal bicolored path connecting two opposite sides, w.l.o.g. top and bottom sides, of the grid. 
We label the columns of the grid as $C_1,\dots,$ $C_{k-2}$ from left to right. Let $C_j$ be the leftmost column containing vertices from $P$. 
If they exist, we label the vertex set in $C_{j-1}, C_j, C_{j+1}, C_{j+2}$ as $v_s, w_s, x_s, y_s,$ $1\leq s\leq k-2$, respectively, $s$ indicating the row index with $s=1$ being the top row. Note that, it is possible that some of these columns do not exist depending on whether $P$ has edges on the right or left side of the grid. We present below a case analysis considering these possibilities as well and find a bicolored $P_k$ in each case. There can be at most one adjacent pair of vertices from different columns (jump across columns) in {\it P}, otherwise we have a bicolored $P_k.$ Thus, $P$ has vertices in at most two columns. Assume, w.l.o.g., that $P$ is a red-blue colored path and $w_1$ has color red in the rest of the proof.  
 

\noindent
{\it Case 1.} $|V(P)|=k-2.$ 

The only possibility of this case is that $P$ contains only vertices from a single column (only $C_j$), i.e. $V(P)=\{w_1,w_2, \dots, w_{k-2}\}.$ Assuming that $C_{j+1}$ exists, $x_1$ has color green by the maximality of $P$. 
Since $w_2$ has color blue, $x_2$ has color red for a proper coloring. Also, the remaining vertices on $C_{j+1}$ have only red and green colors by maximality of $P$. Hence, the vertices on $C_{j+1}$ together with $w_1$ induce a red-green colored $P_{k-1}$. If $C_{j+1}$ does not exist and $P$ is the right side of the grid, then the same arguments yield a red-green colored $P_{k-1}$ similarly using the vertices on $C_{j-1}$ and $w_1.$ This contradicts the maximality of $P.$ 

\noindent
{\it Case 2.} $|V(P)|=k-1.$ 

By symmetry, we assume w.l.o.g. that the jump of $P$ (across columns) is from left to right, specifically from $C_j$ to $C_{j+1}$ with $j\geq 2$ and at the $i$'th row with $2\leq i\leq k-2.$ Thus, $P$ is $w_1, \dots, w_i, x_i, \dots, x_{k-2}.$ 
Note that $C_{j-1}$ exists since $j\geq 2$, and that $w_{i-1}$ and $x_{i-1}$ exist since $i\geq 2.$ By arguments as in Case 1, $v_1$ is (colored) green, $v_2$ is red, $v_3$ is green,..., $v_i$ is red or green (depending on the parity of $i$). The following can be proved similarly. 

\begin{claim}\label{rem:nbrs}
$v_1,\dots,v_i$ and $x_1,\dots,x_{i-1}$ are green-red colored paths with $v_1$ and $x_1$ colored green. If $\alpha$ is the color of $x_{k-2}$, then $w_{k-2},\dots,w_{i+1}$ and $y_{k-2},\dots,y_i$ (if it exists) are green-$\alpha$ colored paths with $w_{k-2}$ and $y_{k-2}$ colored green. 
\end{claim} 
 

If $k$ is even and $i=k-2$, then $v_1,w_1,x_1,\dots,x_{k-2}$ is a green-red colored path by Claim~\ref{rem:nbrs}, contradicting the maximality of $P.$ If $k$ is even, $i<k-2$ and $i$ is even, then $w_{i+1}$ is colored red, which means $w_1,\dots,w_{i+1},x_{i+1},\dots,x_{k-2}$ is also a red-blue colored path. Thus, if $k$ is even and $i<k-2$, we may assume that $i$ is odd. If $k$ is even, $i$ is odd, and $i<k-2$, then $w_1,v_1,\dots,v_i,w_i,\dots,w_{k-2}$ is a red-green colored path by Claim~\ref{rem:nbrs}, contradicting the maximality of $P.$ Thus, we have the following. 
\begin{claim}\label{clm:k_odd}
$k$ is odd.
\end{claim}
If $i$ is odd (i.e., $k$ and $i$ are both odd), then by Claim~\ref{rem:nbrs}, $x_{i-1}$ is colored red and $w_{i+1}$ (if exists) is colored blue, and therefore the path $w_1\dots,w_{i-1},x_{i-1},x_i,w_i,w_{i+1},x_{i+1},\dots,x_{k-2}$ (or the part of it that exists) is a red-blue colored path, contradicting the maximality of $P.$ Thus, we have the following. 
\begin{claim}\label{clm:i_even}
    $i$ is even (which also means that $i<k-2$).
\end{claim}
To complete the proof, it suffices to produce a contradiction assuming that $k$ is odd and $i$ is even. 
If $j\leq k-4$, then $C_{j+2}$ exists and by Claim~\ref{rem:nbrs}, the path $x_i, \dots, x_1, w_1$, $v_1,\dots,v_i$ is a red-green colored $P_{2i+1}$ and 
the path $w_i, \dots, w_{k-2}, x_{k-2}, y_i, \dots, y_{k-2}$ 
is a blue-green colored $P_{2(k-i-1)+1}$. One of these bicolored paths has at least $k$ vertices, depending on the value of $i$.

If $j=k-3$, then $C_{j+2}$ does not exist. In this case, it is not guaranteed to have a bicolored $P_k$ only using vertices in $P$ and their neighbors.  So, we use Lemma~\ref{lem-angle} to show the existence of a bicolored $P_k$.  Let $C$ be the bicolored red-blue component containing $P.$ If $C$ contains any vertex from the top or bottom sides outside $P,$ then $C$ contains a red-blue $P_k.$ If $C$ contains any vertex from the left side, then the red-blue path in $C$ connecting $P$ to the left side together with $P$ yield a bicolored path with at least $k$ vertices. Thus, we assume that $C$ is not connected to any side of the grid (other than the right side) outside $P$. This implies that there exists a partial walk of $C$ from $x_{k-2}$ to $w_1$. However, (4) of Lemma~\ref{lem-angle} yields that $x_{k-2}$ and $w_1$ have the same color, a contradiction. 
\end{proof}

\end{document}
